\documentclass[10pt,a4paper]{amsart}
\usepackage[margin=19mm]{geometry}
\usepackage{amsmath,amssymb,mathtools}
\usepackage[scr=rsfs,cal=euler]{mathalfa}
\usepackage{xfrac}
\usepackage[unicode,hidelinks,bookmarks=false]{hyperref}
\newcommand{\ie}{i.e.\ }
\newcommand{\cf}{cf.\ }
\newcommand{\resp}{respectively\ }
\newcommand{\Subsection}{\subsection}
\providecommand{\nobreakdash}{\nobreak\hbox{-}}
\setlength{\emergencystretch}{2em}
\multlinegap0pt
\usepackage{amssymb}
\usepackage{tikz}
\usepackage[shortlabels]{enumitem}
\newtheorem{theorem}{Theorem}
\usepackage{cleveref}
\crefname{theorem}{Theorem}{Theorems}
\newcommand{\bs}[1]{\ensuremath{\boldsymbol{#1}}}
\DeclareMathOperator{\sym}{sym}
\title{Transform before linearizing: robust Newton methods for singular $p$-Laplace and $p$-Stokes equations}
\author{Gonzalo G. De Diego, Fortino Garcia, Georg Stadler, Kewei Wang}
\date{Source: arXiv:2609.12241v1 (CC BY 4.0)}
\begin{document}
\maketitle

\noindent\emph{Source and licence.} Transform before linearizing: robust Newton methods for singular $p$-Laplace and $p$-Stokes equations. Version arXiv:2609.12241v1 (CC BY 4.0), distributed by arXiv under the Creative Commons Attribution 4.0 International licence, http://creativecommons.org/licenses/by/4.0/, as recorded in the arXiv licence metadata for this exact version and shown on the abstract page https://arxiv.org/abs/2609.12241v1. Full licence notice: \texttt{licenses/arxiv-2609.12241-CC-BY-4.0.txt}.\par\smallskip

\noindent\textit{Editorial scope.} Counted unit: the article's theorem thm:finite-dimensional-lifted-global-convergence (source0.tex lines 530--542). The article's complete original proof of it is delivered with it (source0.tex lines 543--623, one \texttt{proof} environment of 81 lines). No mathematics is written, repaired or re-derived: the statement and the proof are the article's own lines, in the article's own notation, and no retained line is substituted or re-flowed.  Retained uncounted as the source-local context the proof needs: lines 201--207; lines 211--215; lines 313--319; lines 346--348; lines 355--360; lines 366--383; lines 425--427; lines 434--438; lines 444--448; lines 454--456; lines 462--482; lines 489--491; lines 499--501.  Nothing was omitted from any retained span, so the package carries no omission record.  No retained line contains a citation, so the wrapper carries no bibliography and no bibliography entry.  No retained line contains a figure environment, an image-inclusion command or drawing code, so no image is delivered and the package declares no asset dependency.  Editorial formatting only, in addition: an AMS \texttt{amsart} preamble that restates the article's own declarations and macros used by the retained lines, namely the environments \texttt{theorem} on separate counters, together with the standard packages those lines need.  The retained environments print in this document as Theorem 1 in order of appearance, whereas the article prints the same items with the numbering of its own full document.  The complete native source of the article as distributed in the arXiv e-print for this exact version is bundled unchanged as \texttt{sources/210091/source0.tex}.
\begin{equation}
\label{eq:newton-intro}
-\nabla\cdot\left(|\nabla u_k|^{p-2}
\bigl(I + (p-2)\,\bs{w}_k\bs{w}_k^\top\bigr)\nabla\hat u\right)
= \nabla\cdot\bigl(|\nabla u_k|^{p-2}\nabla u_k\bigr) + f
\quad\text{in }\Omega,
\end{equation}

\begin{equation}
\label{eq:picard-intro}
-\nabla\cdot\left(|\nabla u_k|^{p-2}\nabla \hat u\right) = \nabla\cdot\bigl(|\nabla u_k|^{p-2}\nabla u_k\bigr) + f
\quad\text{in }\Omega,
\end{equation}

To isolate the nonlinearity, we thus introduce the $\mathbb R^d$-valued variable  $\bs\lambda(\bs x):= |\nabla u(\bs x)|^{p-2}\, \nabla u(\bs x) $ for all $\bs x\in\Omega$.
The resulting system in the independent variables $(u,\bs\lambda)$ is now
\begin{alignat}{2}
-\nabla\!\cdot\!\bs\lambda &= f \quad&&\text{in }\Omega,\label{eq:extrem1}\\
\bs\lambda - |\nabla u|^{p-2}\nabla u &= 0 \quad&&\text{in }\Omega,\label{eq:extrem2}
\end{alignat}
where we emphasize that \eqref{eq:extrem1} is linear, and the nonlinearity is restricted to the pointwise equation \eqref{eq:extrem2}, which can be manipulated to potentially change the nature of the nonlinearity. Before proceeding with a more detailed discussion of the solution to the lifted system \eqref{eq:extrem1}, \eqref{eq:extrem2}, we rigorously re-derive it below using Fenchel duality theory. We emphasize, however, that it is not necessary for the lifting variable $\bs\lambda$ to admit an interpretation as the solution of a dual problem.

\begin{equation}\label{eq:J-eps}
  J_\varepsilon(u):=\frac1p\int_\Omega|\nabla u|_\varepsilon^{p}\,\mathrm{d}\bs x-\int_\Omega f\,u\,\mathrm{d}\bs x,
\end{equation}

\begin{equation}\label{eq:residual-is-derivative}
  \langle r,v\rangle
  =\int_\Omega |\nabla u|_\varepsilon^{p-2}\nabla u\cdot\nabla v\,\mathrm{d}\bs x
   -\int_\Omega f\,v\,\mathrm{d}\bs x
  = J_\varepsilon'(u)v,
\end{equation}

Introducing an independent variable $\bs\lambda$ as above results in
the following regularized counterparts of \eqref{eq:extrem1}, \eqref{eq:extrem2}:
\begin{alignat}{2}
-\nabla\!\cdot\!\bs\lambda &= f \quad &&\text{in }\Omega,\label{eq:extrem1-reg}\\
\bs\lambda - |\nabla u|_\varepsilon^{p-2}\nabla u &= 0 &&\text{in }\Omega.\label{eq:extrem2-reg}
\end{alignat}
Here, the regularized flux \eqref{eq:extrem2-reg} is defined directly rather than through a dual problem; as noted above, such an interpretation is not required for the lifting.
Since $|\nabla u|_\varepsilon\geq\varepsilon>0$ for all $\bs x\in\Omega$, \eqref{eq:extrem2-reg} is equivalent to
\begin{equation}\label{eq:extrem2-mod-reg}
|\nabla u|_\varepsilon^{2-p}\bs\lambda - \nabla u=0.
\end{equation}
We introduce the corresponding normalized directions
\begin{equation}\label{eq:w-directions}
  \bs w_\varepsilon := \frac{\nabla u}{|\nabla u|_\varepsilon},
  \qquad\qquad
  \bs w_{\bs\lambda} := \frac{\bs\lambda}{|\nabla u|_\varepsilon^{\,p-1}},
\end{equation}
and refer to $\bs w_\varepsilon$ as the \emph{primal direction} and to $\bs w_{\bs\lambda}$ as the \emph{dual direction}. The former is the regularized counterpart of the unit gradient direction $\bs w_k$ used in \eqref{eq:newton-intro}; throughout this section, we suppress the iteration index. Note that $\bs w_{\bs\lambda}=\bs w_\varepsilon$ precisely when $u$ and $\bs\lambda$ satisfy the optimality condition \eqref{eq:extrem2-reg}.

\begin{equation}\label{eq:lifted-tensor}
    \left(I + (p-2)\,\bs w_{\bs\lambda}\bs w_\varepsilon^\top\right)
\end{equation}

\begin{equation}\label{eq:tensor-quadratic-form}
  \bs z^\top\bigl(I + (p-2)\,\bs w_{\bs\lambda}\bs w_\varepsilon^\top\bigr)\bs z
  = |\bs z|^2 + (p-2)\,(\bs w_{\bs\lambda}\cdot\bs z)(\bs w_\varepsilon \cdot \bs z)
  \;\ge\; \bigl(1-(2-p)\,|\bs w_{\bs\lambda}|\,|\bs w_\varepsilon|\bigr)|\bs z|^2 .
\end{equation}

\begin{equation}\label{eq:tensor-sandwich}
  (1-\beta)\,|\bs z|^2
  \;\le\; \bs z^\top\bigl(I + (p-2)\,\bs w_{\bs\lambda}\bs w_\varepsilon^\top\bigr)\bs z
  \;\le\; (1+\beta)\,|\bs z|^2 ,
\end{equation}

\begin{equation}\label{eq:beta-window}
  2-p<\beta<1.
\end{equation}

\begin{equation}\label{eq:lifted-tensor-sym}
  I + (p-2)\,\sym\bigl(\bs w_{\bs\lambda}\bs w_\varepsilon^\top\bigr),
  \qquad \sym(B):=\frac12\bigl(B+B^\top\bigr).
\end{equation}
This leaves the associated quadratic form---and hence the bounds \eqref{eq:tensor-quadratic-form} and \eqref{eq:tensor-sandwich} established above---unchanged. It does, however, render the discretized operator symmetric, allowing the use of symmetric solvers and preconditioners for the resulting linear systems. Observe that this modification has no effect at a consistent pair, where $\bs w_{\bs\lambda}=\bs w_\varepsilon$ and \eqref{eq:lifted-tensor} is already symmetric.

Note that \eqref{eq:tensor-sandwich} is a statement about well-posedness, not about conditioning: since the lifted and the standard linearizations coincide at the solution, the two methods cannot differ in their asymptotic conditioning. What \eqref{eq:tensor-sandwich} does provide is a linear system that remains solvable, and an increment that remains a descent direction for $J_\varepsilon$, no matter how far the dual direction lags behind the primal one. This is what underlies the global convergence result of the next section.


\subsection{Summary of lifted Newton algorithm}
\label{ssec:lifted-algorithm}
Collecting the ingredients from the previous subsections yields the method we analyze and test below. We write
$\bs w_{\varepsilon,k}:=\nabla u_k/|\nabla u_k|_\varepsilon$ and
$\bs w_{\bs\lambda,k}:=\bs\lambda_k/|\nabla u_k|_\varepsilon^{\,p-1}$
for the directions \eqref{eq:w-directions} evaluated at the $k$-th iterate, and recall the regularized objective \eqref{eq:J-eps} with corresponding (lifted) stationarity conditions \eqref{eq:extrem1-reg}, \eqref{eq:extrem2-reg}.

We assume that $\beta$ satisfies \eqref{eq:beta-window}, and enforce the feasibility bound only at $\bs x\in\Omega$ where it is needed, namely in the construction of the diffusion tensor. We define
\begin{equation}\label{eq:discrete-lambda-projection}
  \bs\lambda_k^{\mathrm{pr}}(\bs x):=\frac{\bs\lambda_k(\bs x)}{\max\bigl(1,\,|\bs\lambda_k(\bs x)|/\rho_k(\bs x)\bigr)},
  \qquad \rho_k(\bs x):=\frac{\beta}{2-p}\,|\nabla u_k(\bs x)|_\varepsilon^{p-1},
\end{equation}

\begin{equation}\label{eq:Mk}
  M_k:=|\nabla u_k|_\varepsilon^{p-2}\Bigl(I+(p-2)\,\sym\bigl(\bs w_{\bs\lambda,k}^{\mathrm{pr}}\bs w_{\varepsilon,k}^\top\bigr)\Bigr),
\end{equation}

\begin{equation}\label{eq:armijo}
  J_\varepsilon(u_k+\alpha_k\hat u_k)\le J_\varepsilon(u_k)+\gamma\,\alpha_k\,J_\varepsilon'(u_k)\hat u_k.
\end{equation}

\begin{theorem}[Global convergence of the lifted method]
\label{thm:finite-dimensional-lifted-global-convergence}
Let $\beta$ satisfy \eqref{eq:beta-window} and assume arbitrary initializations $u_{h,0}\in g+\mathring{V}_{h}$
and 
$\bs\lambda_{h,0}$. 
Then the (discrete) iterates produced by the iteration of \cref{ssec:lifted-algorithm} converge globally, in the sense that the discrete residual
$r_{h,k}\to 0$ in $\mathring{V}_{h}^\star$,
and
\begin{equation*}
  u_{h,k}\to u_h^* \quad\text{in }V_h,
\end{equation*}
where $u_h^*\in g+\mathring{V}_{h}$ is the unique minimizer of $J_\varepsilon$ over $g+\mathring{V}_{h}$.
\end{theorem}

\begin{proof}
Consider the sublevel set
$\mathcal L_h:=\{u_h\in g+\mathring{V}_h:J_\varepsilon(u_h)\le J_\varepsilon(u_0)\}$, which is bounded because
$J_\varepsilon(u_h)\to \infty$ as $\|u_h\|_h\to \infty$. %
Since $V_h$ is finite-dimensional, $\mathcal{L}_h$ is compact, and therefore
$ C_h:=\max_{u_h\in\mathcal L_h}\|\nabla u_h\|_{L^\infty(\Omega)}<\infty$.
We prove the uniform boundedness and coercivity of the linearization resulting from the lifting approach on the sublevel set. To this end, suppose that $u_{h,k}\in\mathcal L_h$.
We first derive pointwise spectral bounds for $M_k$ defined in \eqref{eq:Mk}. For any $\bs z\in\mathbb R^d$,
\eqref{eq:tensor-sandwich} gives
\begin{equation}\label{eq:spectral-lower-upper}
  \begin{aligned}
    \bs z^\top M_k \bs z
    &\geq (1-\beta)|\nabla u_{h,k}|_\varepsilon^{p-2}|\bs z|^2
    \geq \underline m_h
    |\bs z|^2,\\
  \bs z^\top M_k \bs z&\leq (1+\beta)|\nabla u_{h,k}|_\varepsilon^{p-2}|\bs z|^2\leq
  \overline m_h |\bs z|^2,
  \end{aligned}
\end{equation}
where $\underline m_h:=(1-\beta)(C_h^2+\varepsilon^2)^{(p-2)/2}$ and  $\overline m_h:=(1+\beta)\varepsilon^{p-2}$.
Denoting
\begin{equation*}
  a_k(v_h,w_h):=\int_\Omega M_k\nabla v_h\cdot\nabla w_h\,\mathrm{d}\bs x,
\end{equation*}
we thus have
\begin{equation}\label{eq:Mk-coercive-bounded}
  a_k(v_h,v_h)\geq \underline m_h\|v_h\|_h^2, \quad |a_k(v_h,w_h)| \leq \overline m_h\|v_h\|_h\|w_h\|_h \quad\text{ for all } v_h,w_h\in\mathring{V}_h.
\end{equation}
The pointwise bounds \eqref{eq:Mk-coercive-bounded} imply that $a_k$ defines an inner product on $\mathring V_h$, so the primal increment is uniquely defined by $a_k(\hat u_{h,k},v_h)=-J_\varepsilon'(u_{h,k})v_h$ for all $v_h\in\mathring V_h$. Choosing $v_h=\hat u_{h,k}$ gives
\begin{equation}\label{eq:direct-descent-identity}
  J_\varepsilon'(u_{h,k})\hat u_{h,k}
  =-a_k(\hat u_{h,k},\hat u_{h,k}).
\end{equation}
If $\hat u_{h,k}=0$, then $J_\varepsilon'(u_{h,k})=0$ on
$\mathring V_h$, and strict convexity shows that $u_{h,k}=u_h^*$.  We may
therefore assume below that $\hat u_{h,k}\neq0$.

Next, we bound the energy along the search direction. Since $p<2$, the map $t\mapsto t^{p/2}$ is concave and thus lies below its tangent line at any $t_0>0$.
Applying this pointwise with $t_0=|\nabla u_h|_\varepsilon^2$ and $t=|\nabla u_h+\alpha\nabla v_h|_\varepsilon^2$, whose difference is $2\alpha\nabla u_h\!\cdot\!\nabla v_h+\alpha^2|\nabla v_h|^2$, then dividing by $p$, integrating, and adding the linear load term to both sides gives
\begin{equation}\label{eq:energy-majorization-discrete}
  J_\varepsilon(u_h+\alpha v_h)
  \le J_\varepsilon(u_h)+\alpha J_\varepsilon'(u_h)v_h
  +\frac{\alpha^2}{2}\int_\Omega
  |\nabla u_h|_\varepsilon^{p-2}|\nabla v_h|^2\,\mathrm{d}\bs x
\end{equation}
for all $u_h\in g+\mathring V_h$, $v_h\in\mathring V_h$, and $\alpha\ge0$. Note that the last term in \eqref{eq:energy-majorization-discrete} is the bilinear form of the Picard iteration \eqref{eq:picard-intro}, which majorizes the curvature of $J_\varepsilon$ because $p<2$. The estimate \eqref{eq:energy-majorization-discrete} holds globally in $\alpha$ and involves no remainder term.

We now take $u_h=u_{h,k}$ and $v_h=\hat u_{h,k}$. The lower bound in \eqref{eq:tensor-sandwich} gives
\begin{equation*}
  \int_\Omega|\nabla u_{h,k}|_\varepsilon^{p-2}|\nabla\hat u_{h,k}|^2\,\mathrm{d}\bs x
  \le\frac{1}{1-\beta}\,a_k(\hat u_{h,k},\hat u_{h,k}),
\end{equation*}
so that \eqref{eq:energy-majorization-discrete} and \eqref{eq:direct-descent-identity} yield, for every $\alpha\ge0$,
\begin{equation}\label{eq:armijo-majorant}
  J_\varepsilon(u_{h,k}+\alpha\hat u_{h,k})
  \le J_\varepsilon(u_{h,k})
  -\alpha\Bigl(1-\frac{\alpha}{2(1-\beta)}\Bigr)
   a_k(\hat u_{h,k},\hat u_{h,k}).
\end{equation}
By \eqref{eq:direct-descent-identity}, the Armijo condition \eqref{eq:armijo} requires a decrease of at least $\gamma\alpha\,a_k(\hat u_{h,k},\hat u_{h,k})$. Comparing with \eqref{eq:armijo-majorant}, it is therefore satisfied as soon as the bracket is at least $\gamma$, that is,
whenever $0<\alpha\le2(1-\gamma)(1-\beta)$. Since the line search halves the trial step
starting from $\alpha=1$, it terminates with an accepted step size
$\alpha_k\ge\underline\alpha:=(1-\gamma)(1-\beta)>0$.
Since \eqref{eq:armijo} and \eqref{eq:direct-descent-identity} give $J_\varepsilon(u_{h,k+1})\le J_\varepsilon(u_{h,k})$, starting from
$u_{h,0}\in\mathcal L_h$, induction therefore proves that the iteration is
well defined and that every $u_{h,k}$ remains in $\mathcal L_h$.
Using \eqref{eq:direct-descent-identity}, \eqref{eq:armijo}, and the lower bound on $\alpha$, we obtain
\begin{equation*}
  J_\varepsilon(u_{h,k})-J_\varepsilon(u_{h,k+1})
  \ge \gamma\underline\alpha\,
  a_k(\hat u_{h,k},\hat u_{h,k}).
\end{equation*}
Since $J_\varepsilon$ is bounded below, summation over $k$ shows that $a_k(\hat u_{h,k},\hat u_{h,k})\to 0$. Cauchy--Schwarz in the  $a_k$-inner product and the upper bound in
\eqref{eq:Mk-coercive-bounded} give, for every $v_h\in\mathring V_h$,
\begin{equation*}
  |J_\varepsilon'(u_{h,k})v_h|^2
  =|a_k(\hat u_{h,k},v_h)|^2
  \le \overline m_h\, a_k(\hat u_{h,k},\hat u_{h,k})\|v_h\|_h^2,
\end{equation*}
so that $r_{h,k}\to0$ in $\mathring V_h^\star$ by \eqref{eq:residual-is-derivative}. Compactness of $\mathcal L_h$ and strict convexity of $J_\varepsilon$ then prove $u_{h,k}\to u_h^*$ in $V_h$.
\end{proof}

\end{document}
